\section{Complete integer witnesses for abnormal single gaps}
\label{sec:single-gap}

\subsection{Canonical model and comparison threshold}

For an integer \(g\geq1\), set
\[
 D_g=\{-4j:j\geq0\}\cup\{0,g\}\cup\{g+4j:j\geq0\}.
\]
Let \(Q_i=+1\) for \(i\in D_g\) and \(Q_i=-1\) otherwise.  Choose either
lift \(\tau_{i+1}=Q_i\tau_i\), and define on \(\ell^2(\Z)\)
\begin{equation}
 (A_gv)_k=v_{k-1}+v_{k+1}+\tau_{k-2}v_{k-2}+\tau_kv_{k+2},
 \qquad H_g=A_g^2.
 \label{eq:single-gap-operator}
\end{equation}
Gap four is the unperturbed period-eight bulk.  Gap six is the elementary
interface of \cref{sec:exact-2r}.

For a finitely supported integer vector \(v\), extend it by zero.  Since
\(A_g\) is self-adjoint,
\begin{equation}
 \frac{\langle v,H_gv\rangle}{\langle v,v\rangle}
 =\frac{\|A_gv\|^2}{\|v\|^2}.
 \label{eq:rayleigh}
\end{equation}
The certified upper endpoint in \eqref{eq:c6-interval} shows that it is
enough to exceed the rational threshold
\begin{equation}
 T=\frac{7905369311620328}{10^{15}}+\frac1{250}
 =\frac{988671163952541}{125000000000000}.
 \label{eq:T}
\end{equation}
For integers \(N,D>0\), define the exact cross-multiplication
\begin{equation}
 \Delta(N,D)=125000000000000N-988671163952541D.
 \label{eq:Delta}
\end{equation}
Then \(N/D>T\) if and only if \(\Delta(N,D)>0\).

\subsection{Six small-gap witnesses}

Coordinates in each vector are listed in increasing order on the interval
\(I\).  The displayed image is listed on \(J=I+[-2,2]\).  Every entry follows
directly from the four-term formula \eqref{eq:single-gap-operator}.

\paragraph{Gap \(g=1\).}
On \(I=[-2,3]\), take
\begin{align*}
 v&=(2,0,4,4,6,5),\\
 A_1v&=(-2,2,4,2,12,15,13,10,11,-5).
\end{align*}
Then \(D=\|v\|^2=97\) and \(N=\|A_1v\|^2=812\).

\paragraph{Gap \(g=2\).}
On \(I=[-4,6]\), take
\begin{align*}
 v&=(1,-1,-2,-1,-4,-5,-6,-2,1,4,2),\\
 A_2v&=(-1,0,1,0,-7,0,-14,-11,-12,-14,10,5,5,-2,2).
\end{align*}
Then \(D=109\) and \(N=866\).

\paragraph{Gap \(g=3\).}
On \(I=[-5,8]\), take
\begin{align*}
 v={}&(0,1,-1,-3,0,-6,-8,-6,-10,-6,-8,-6,1,-3),\\
 A_3v={}&(0,-1,0,2,-2,-8,0,-17,-22,-18,-28,-18,-23,-16,\\
 &\hspace{31mm}-1,-5,-4,3).
\end{align*}
Then \(D=393\) and \(N=3114\).

\paragraph{Gap \(g=5\).}
On \(I=[-2,7]\), take
\begin{align*}
 v&=(2,0,4,4,4,4,1,3,3,3),\\
 A_5v&=(-2,2,4,2,10,12,11,12,0,11,5,6,6,-3).
\end{align*}
Then \(D=96\) and \(N=764\).

\paragraph{Gap \(g=7\).}
On \(I=[-2,9]\), take
\begin{align*}
 v&=(2,0,3,4,4,4,1,3,2,3,2,3),\\
 A_7v&=(-2,2,3,1,10,11,10,12,1,10,3,10,4,5,5,-3).
\end{align*}
Then \(D=97\) and \(N=768\).

\paragraph{Gap \(g=8\).}
On \(I=[-8,16]\), take
\begin{align*}
 v={}&(4,4,4,3,-3,3,9,1,19,22,21,22,4,16,8,12,5,6,0,-4,\\
 &\hspace{31mm}-1,1,0,1,0),\\
 A_8v={}&(4,0,8,11,14,8,-7,8,26,3,53,61,59,63,9,46,19,35,\\
 &\hspace{31mm}10,21,-4,-8,-3,4,1,1,1,1,0).
\end{align*}
Then \(D=2487\) and \(N=19672\).

All exact comparisons for these six witnesses are collected in
\cref{tab:small-gap-margins}.  The third column is the raw integer from
\eqref{eq:Delta}; the last column is the reduced or partially reduced exact
difference from \(T\).

\begin{table}[htbp]
\centering
\small
\caption{Exact small-gap Rayleigh comparisons.}
\label{tab:small-gap-margins}
\resizebox{\textwidth}{!}{%
\begin{tabular}{@{}ccll@{}}
\toprule
\(g\)&\(N/D\)&\(\Delta(N,D)\)&\(N/D-T\)\\
\midrule
1 & \(812/97\) & \(5598897096603523\) &
 \(5598897096603523/12125000000000000\)\\
2 & \(866/109\) & \(484843129173031\) &
 \(484843129173031/13625000000000000\)\\
3 & \(3114/393\) & \(702232566651387\) &
 \(234077522217129/16375000000000000\)\\
5 & \(764/96\) & \(587568260556064\) &
 \(18361508142377/375000000000000\)\\
7 & \(768/97\) & \(98897096603523\) &
 \(98897096603523/12125000000000000\)\\
8 & \(19672/2487\) & \(174815250030533\) &
 \(174815250030533/310875000000000000\)\\
\bottomrule
\end{tabular}
}
\end{table}

Every integer and every rational number in \cref{tab:small-gap-margins} is
positive.  Thus \eqref{eq:rayleigh} proves
\(\sup\spec(H_g)>\cSix+1/250\) for these six gaps.

\subsection{One witness for every \texorpdfstring{\(g\geq9\)}{g at least 9}}

On \(I=[-2,11]\), use the fixed vector
\begin{equation}
 v_*=(4,0,7,8,8,9,1,7,3,6,1,4,1,2),
 \qquad \|v_*\|^2=391.
 \label{eq:vstar}
\end{equation}
Only coefficients \(\tau_i\) with \(-4\leq i\leq11\) enter its nonzero
image.  The right defect can therefore affect the calculation only for
\(g=9\) or \(g=10\); all \(g\geq11\) have the same local coefficient word.
On \(J=[-4,13]\), direct evaluation gives
\begin{align}
 A_9v_*={}&(-4,4,7,3,20,24,23,24,5,19,11,15,6,10,5,5,3,-2),
 \label{eq:g9-image}\\
 A_{10}v_*={}&(-4,4,7,3,20,24,23,24,5,19,11,15,6,10,5,5,1,-2),
 \label{eq:g10-image}\\
 A_gv_*={}&(-4,4,7,3,20,24,23,24,5,19,11,15,6,10,5,5,1,2),
 \quad(g\geq11). \label{eq:gt10-image}
\end{align}
Moreover,
\[
\|A_9v_*\|^2=3102,\qquad
\|A_{10}v_*\|^2=3094,\qquad
\|A_gv_*\|^2=3094\quad(g\geq11).
\]
Hence every \(g\geq9\) satisfies
\begin{equation}
 \sup\spec(H_g)\geq\frac{3094}{391}=\frac{182}{23}.
\end{equation}
The final exact comparison is
\begin{equation}
 \Delta(182,23)=10563229091557>0,
 \qquad
 \frac{182}{23}-T=
 \frac{10563229091557}{2875000000000000}>0.
 \label{eq:tail-margin}
\end{equation}
Finite propagation, rather than sampling, is what reduces the infinite tail
of gap lengths to the three images \eqref{eq:g9-image}--\eqref{eq:gt10-image}.

\subsection{Completion of the single-gap hierarchy}

For \(g=4\), the defect lattice is \(4\Z\), so the operator is the
period-eight reference bulk and
\begin{equation}
 \sup\spec(H_4)=\etaRef<\cSix.
\end{equation}
For \(g=6\), the global interface theorem gives
\begin{equation}
 \sup\spec(H_6)=\cSix,\qquad
 \dim\ker(H_6-\cSix)=2.
\end{equation}
The six small vectors and the fixed vector \(v_*\) exhaust every positive
gap outside \(\{4,6\}\).  Equation \eqref{eq:lift-equivalence} transfers the
conclusion between the two lifts, and reflection transfers it between the
two orientations.  The smallest printed margin is the \(g=8\) value
\[
 \frac{174815250030533}{310875000000000000}>0.
\]
We have proved the following strict statement.

\begin{theorem}[Uniform abnormal single-gap separation]
\label{thm:single-gap}
For either lift and either interface orientation,
\[
 \sup\spec(H_6)=\cSix,\qquad
 \dim\ker(H_6-\cSix)=2.
\]
For every positive integer \(g\notin\{4,6\}\),
\[
 \sup\spec(H_g)>\cSix+\frac1{250}.
\]
Thus G6 is the unique minimizer among abnormal positive single gaps.  The
statement does not compare G6 with arbitrary multi-gap or finite-core
interfaces.
\end{theorem}
